\section{Global equivariant classes}
\label{sec:equivariant}

We now compute the global reflection cohomology and construct the
full-group classes needed in the doubling argument.
The computation uses a round metric only to filter the loop space;
the resulting classes belong to the same global groups for every
metric.  The distinction between odd and even round stages will be
essential.  It is independent of the parity of the sphere dimension.

\subsection{Natural Gysin sequences}

Let \(w\in H^1(BG;\kk)\) and \(U\in H^2(BG;\kk)\) be the
Stiefel--Whitney classes of the standard two-dimensional representation
of \(G=O(2)\).  We use the same letter \(w\) for its restrictions.
Let \(v\in H^1(BK;\kk)\) be the character which is nontrivial on
\(\tau\) and trivial on \(R\).  Then
\begin{equation}
 H^*(BG;\kk)=\kk[w,U],\qquad
 H^*(BK;\kk)=\kk[v,w],\qquad
 \res_K U=v(v+w).
 \label{eqv:coefficient-rings}
\end{equation}
Indeed, the standard plane restricts to the two characters \(v\) and
\(v+w\).  Its projectivization over \(BG\) is \(BK\), so the
projective-bundle formula makes restriction injective
\cite[Section~3.1, pp.~78--79]{HatcherVBKT2017}.  Conjugation
in \(G\) interchanges the two coordinate lines of \(K\), whereas inner
conjugation acts trivially on \(H^*(BG)\).  The image is consequently
contained in the symmetric polynomials in \(v,v+w\).  It contains their
sum \(w\) and product \(U\), which generate that symmetric polynomial
ring.  This proves \eqref{eqv:coefficient-rings}.

For an \(H\)-pair \((Y,A)\), write \(H_H^*(Y,A)\) for its Borel
cohomology.  These groups are zero in negative degrees.

\begin{proposition}[Natural Gysin sequences]
\label{prop:gysin}
For every \(G\)-pair \((Y,A)\), there are natural exact sequences
\begin{align}
 \cdots\longrightarrow H_G^{j-2}(Y,A)
 &\xrightarrow{\,U\,} H_G^j(Y,A)
 \xrightarrow{\res_R} H_R^j(Y,A)
 \xrightarrow{\int} H_G^{j-1}(Y,A)
 \xrightarrow{\,U\,} H_G^{j+1}(Y,A)\longrightarrow\cdots,
 \label{eqv:standard-gysin}\\
 0\longrightarrow H_G^j(Y,A)
 &\xrightarrow{\res_K}H_K^j(Y,A)
 \xrightarrow{\int}H_G^{j-1}(Y,A)
 \longrightarrow0.
 \label{eqv:weight-two-gysin}
\end{align}
For an \(R\)-pair and a \(K\)-pair, respectively, the sign-line
sequences are
\begin{align}
 \cdots\longrightarrow H_R^{j-1}(Y,A)
 &\xrightarrow{\,w\,}H_R^j(Y,A)
 \longrightarrow H^j(Y,A)
 \longrightarrow H_R^j(Y,A)
 \xrightarrow{\,w\,}H_R^{j+1}(Y,A)\longrightarrow\cdots,
 \label{eqv:reflection-gysin}\\
 \cdots\longrightarrow H_K^{j-1}(Y,A)
 &\xrightarrow{\,v\,}H_K^j(Y,A)
 \longrightarrow H_R^j(Y,A)
 \longrightarrow H_K^j(Y,A)
 \xrightarrow{\,v\,}H_K^{j+1}(Y,A)\longrightarrow\cdots.
 \label{eqv:half-shift-gysin}
\end{align}
The maps commute with maps of pairs and subgroup restriction wherever
the latter applies, and satisfy the projection formula for absolute
base classes.
\end{proposition}

\begin{proof}
For a real rank-\(r\) representation of \(H\), its associated bundle
over \(Y_H\) restricts to a bundle over \(A_H\).
Suppose that its unit sphere is \(H/H'\).
The sphere-bundle pair is then the Borel pair for \(H'\).
The relative Thom sequence gives
\begin{equation}
 \cdots\longrightarrow H_H^{j-r}(Y,A)
 \xrightarrow{\,w_r\,}H_H^j(Y,A)
 \longrightarrow H_{H'}^j(Y,A)
 \xrightarrow{\int}H_H^{j-r+1}(Y,A)\longrightarrow\cdots.
 \label{eqv:general-gysin}
\end{equation}

Here is a cochain construction that also justifies
\eqref{eqv:general-gysin} when the inclusion \(A_H\subset Y_H\) is not
a cofibration.  For a vector bundle over \(Z\supset A_0\), let \(D,S\)
denote its disk and sphere bundles, and let \(D_{A_0},S_{A_0}\)
be their restrictions.  In the singular chain complex of \(D\), put
\[
 \mathcal L=C_*(S)+C_*(D_{A_0}).
\]
The two summands intersect in \(C_*(S_{A_0})\), since singular chains
have the singular simplices as a basis.  Restriction of cochains gives
degreewise exact sequences
\begin{align*}
 0&\longrightarrow\operatorname{Ann}\mathcal L
 \longrightarrow C^*(D,S)
 \longrightarrow C^*(D_{A_0},S_{A_0})\longrightarrow0,\\
 0&\longrightarrow\operatorname{Ann}\mathcal L
 \longrightarrow C^*(D,D_{A_0})
 \longrightarrow C^*(S,S_{A_0})\longrightarrow0.
\end{align*}
Surjectivity follows by extending linear functionals on chain
subspaces.  Cup with a Thom cocycle sends a relative base cochain into
\(\operatorname{Ann}\mathcal L\): on \(S\) the Thom factor vanishes,
and on \(D_{A_0}\) the base factor vanishes.
Comparison of the first sequence with the base-pair sequence, using
the absolute Thom isomorphisms, proves
\[
 H^{j-r}(Z,A_0)\cong H^j(\operatorname{Ann}\mathcal L).
\]
The second sequence and the disk contraction give
\eqref{eqv:general-gysin}.  The map back to the base is cup with the
zero-section pullback of the Thom class, namely \(w_r\).
Pullback of Thom classes gives naturality and the projection formula.
The absolute Thom isomorphism used here is the ordinary mod-two
one \cite[Section~3.2, pp.~88--89]{HatcherVBKT2017};
for a general base it follows by pulling back to its singular
CW replacement and using the disk and sphere fibrations.
This argument uses the sum of two chain subspaces and does not require
an excision assertion for their union.

The standard \(G\)-plane has sphere stabilizer \(R\) and Euler class
\(U\), proving \eqref{eqv:standard-gysin}.
Let \(\rho_2\) be the representation that doubles rotation angles and
sends the chosen reflection to \(\operatorname{diag}(1,-1)\).
Its sphere stabilizer is \(K\).  On \(K\) it has a fixed line, so its
Euler class restricts to zero.  The injectivity in
\eqref{eqv:coefficient-rings} makes that Euler class zero already on
\(BG\), giving \eqref{eqv:weight-two-gysin}.
Finally, the sign representations with kernels \(1\subset R\) and
\(R\subset K\) have Euler classes \(w\) and \(v\), respectively.
\end{proof}

\subsection{The classes to be constructed}

Evaluation at zero is \(R\)-invariant.  If
\(\sigma\in H^n(S^n;\kk)\) is the nonzero class, set
\[
 e=\ev_0^*\sigma\in H_R^n(\Lambda;\kk).
\]
In particular \(e^2=0\).  Recall that \(d=n-1\), and put
\[
 a_m=(2m-1)d\qquad(m\ge1).
\]

\begin{theorem}[Global generators and full-group lifts]
\label{thm:global-generators}
For \(n\ge3\), there are classes
\[
 X_m\in B_\infty^{a_m}=H_R^{a_m}(\Lambda,M),\qquad m\ge1,
\]
with the following properties.
\begin{enumerate}[label=\textup{(\roman*)}]
\item The classes
\[
 \xi_{m,k}=w^kX_m,\qquad
 \eta_{m,k}=e\,w^kX_m,
 \qquad m\ge1,\quad 0\le k\le d,
\]
form a graded \(\kk\)-basis of \(B_\infty\).  Their degrees are
\[
 |\xi_{m,k}|=a_m+k,\qquad
 |\eta_{m,k}|=a_m+n+k.
\]
We call the \(\xi_{m,k}\) the \emph{bottoms} and the \(\eta_{m,k}\)
the \emph{tops} at stage \(m\).
Multiplication by \(w^n\) annihilates \(B_\infty\).
\item Every \(X_m\) is the restriction of a class
\(\zeta_m\in H_G^{a_m}(\Lambda,M)\).
If \(m\) is odd, then \(U^n\zeta_m=0\).
If \(m\) is even, there is also a class
\(Y_m\in H_G^{a_m+n}(\Lambda,M)\) with
\(\res_R Y_m=eX_m\).
\item For the round metric, let \(Z_p\) be an open regular sublevel
containing exactly the first \(p\) positive critical manifolds,
and choose \(Z_0\) below the first one.
For \(H=R,G\), every crossing sequence for
\((Z_p,Z_{m_0-1})\), \(p\ge m_0\), is short exact.
A chosen class of the \(m\)-th crossing has a unique extension to
\(H_H^j(\Lambda,Z_{m-1})\) whenever \(j<a_{m+1}\), for
\(H=R,G\).  The lifts in \textup{(ii)} and the classes \(X_m\)
can be obtained from these relative extensions.
\end{enumerate}
\end{theorem}

The top classes in \textup{(ii)} are chosen lifts; they need not be
canonical.  Only their reflection restrictions will be used.
Also, the annihilation statement in \textup{(i)} is relative to
constants.  The absolute constant-loop Borel summand retains its
polynomial \(w\)-tower.

\subsection{Round crossings}

Give \(S^n\) its unit round metric and write \(E_0\) for its energy.
The positive critical values are \(c_m=(2\pi m)^2\).
The critical manifold at \(c_m\) consists of the loops
\begin{equation}
 \gamma_{x,u}(t)=x\cos(2\pi mt)+u\sin(2\pi mt),
 \qquad |x|=|u|=1,\quad x\perp u.
 \label{eqv:round-loops}
\end{equation}
It is \(C_m=STS^n\), the manifold of orthonormal two-frames in
\(\mathbb R^{n+1}\).

\begin{lemma}[Round attachment]
\label{eqv:round-attachment}
The negative bundle of \(C_m\) has rank \(a_m\), and the kernel of the
round Hessian is exactly \(TC_m\).
For every closed subgroup \(H\subset G\), the crossing has a Thom
isomorphism
\begin{equation}
 H_H^j(Z_m,Z_{m-1})\cong H_H^{j-a_m}(C_m).
 \label{eqv:crossing-thom}
\end{equation}
These isomorphisms commute with subgroup restriction and cup by
absolute classes.  The small sublevel \(Z_0\) retracts
\(G\)-equivariantly to \(M\).
\end{lemma}

\begin{proof}
Along \eqref{eqv:round-loops}, put
\[
 e_t=-x\sin(2\pi mt)+u\cos(2\pi mt).
\]
A tangent field is \(a(t)e_t+b(t)\), where
\(b(t)\perp\operatorname{span}\{x,u\}\).
Half the Hessian is
\begin{equation}
 \int_0^1\bigl(|a'|^2+|b'|^2-(2\pi m)^2|b|^2\bigr)\,dt.
 \label{eqv:round-hessian}
\end{equation}
In each of the \(d\) normal directions, the negative Fourier
frequencies are \(|k|<m\).  Their real dimension is \(2m-1\).
The null modes are the two normal modes at frequency \(m\), together
with the constant tangential mode.  They have dimension \(2d+1\)
and are precisely the variations of the frame \((x,u)\).
The nonzero Fourier eigenvalues, divided by the corresponding squared
\(H^1\)-norms, are bounded away from zero.  Thus the Hessian is
invertible on the orthogonal complement of its kernel.
At constants, the Hessian is \(2\int|V'|^2\), positive and invertible
off its constant kernel.

We describe why the resulting attachment respects the full group.
Use the strong metric induced by the Euclidean \(H^1\)-norm.
Each parameter symmetry is a smooth isometry, and the action is
continuous.  Pointwise exponential gives a chart on the Hilbert
normal bundle of the compact manifold \(C_m\), at a uniform small
radius.  In these coordinates Taylor's formula writes the energy
difference as \(\langle A(z,y)y,y\rangle\), with
\(A(z,0)\) an invertible symmetric operator on the normal fiber.
Normalize first by \(|A(z,0)|^{1/2}\), and let
\(J=\operatorname{sign}A(z,0)\).
For the resulting operator \(B\) near \(J\), the norm-convergent
power series \(C=(JB)^{1/2}\) satisfies \(C^*JC=B\).
The change \(y\mapsto C(z,y)y\) is a diffeomorphism after shrinking
the tube.  All these constructions commute with the parameter
isometries.  They give the exact normal form
\begin{equation}
 E_0=c_m-|u_-|^2+|u_+|^2
 \label{eqv:round-normal-form}
\end{equation}
on the negative and positive bundles.

For completeness, the sublevel attachment uses the original Hilbert
space, rather than a polygon model with only finite rotational
symmetry.  Patch a complete normalized descending gradient field to
the split field
\[
 \dot z=0,\qquad \dot u_-=u_-,\qquad \dot u_+=-u_+
\]
inside the normal tube, using an invariant squared-radius cutoff.
Both fields strictly decrease \(E_0\) away from \(C_m\) within the
isolating energy band.
The coordinate derivatives are uniformly bounded in a small tube over
the compact zero section, so the patched field remains bounded.
The ambient strong metric is complete; hence this field is complete.
Palais--Smale and regular-band deformation
(Lemmas~\ref{lem:local-analytic} and~\ref{lem:local-strips})
give a uniform decrease outside any smaller tube in an isolating
energy band.
The flow-lens construction of Lemma~\ref{lem:local-lens}
therefore computes the crossing in that smaller tube.
Use its notation \(X,Y,N,Q\), with lower energy cut \(t_-<c_m\).
In split coordinates, contraction of \(u_+\) decreases both endpoint
energies of a trapped flow segment.  It defines maps \(N\to X\)
taking \(Q\) into \(Y\), even when a contracted point leaves \(N\)
through the lower endpoint.  Via \(N/Q\cong X/Y\), these maps induce
a contraction to the negative disk with its boundary collapsed.
For flow time \(T\), the surviving disk before applying the flow has
radius \(e^{-T}\sqrt{c_m-t_-}\); its boundary flows to the lower cut.
The flow multiplies \(u_-\) by \(e^T\), identifying this disk and
sphere with the usual negative-bundle disk and sphere.
The disk is closed in the lens and meets \(Q\) exactly in its sphere.
These statements remain true after Borel construction.
The disk/sphere quotient consequently embeds in the lens quotient:
the saturation of a closed subset of the disk is that subset,
possibly together with the closed collapse set.
The contraction fixes its image.
The disk/sphere inclusion is a good pair.  This proves
\eqref{eqv:crossing-thom} by the mod-two Thom isomorphism.

All pair maps just used are \(G\)-equivariant.  The negative bundle
carries its actual derivative action; no triviality of that action
is assumed.  Restriction to a subgroup and cup by an absolute class
therefore commute with the attachment isomorphism.
Wider regular cuts give the same assertion by regular-band
deformation.  Finally, sufficiently small energy places a loop in
the positive normal chart of the constants, where contraction of
the normal coordinate proves the last assertion.
\end{proof}

\subsection{Cohomology of the critical manifolds}

Reflection sends \((x,u)\) to \((x,-u)\).
It acts freely on \(C_m\), and
\[
 (C_m)_R\simeq P(TS^n).
\]
The sign line is the tautological line in this projective bundle.
Since \(TS^n\oplus\mathbb R\) is trivial, its positive
Stiefel--Whitney classes vanish.  The projective-bundle formula gives
\begin{equation}
 H_R^*(C_m)=\kk[w,e]/(w^n,e^2).
 \label{eqv:reflection-critical-ring}
\end{equation}
Here \(e\) is pulled back from the first frame vector \(x\).
This computation uses mod-two coefficients also when \(n\) is even.
Write \(x_m\) for the bottom Thom class of the crossing, with the
group understood.  Its subgroup restrictions are the corresponding
bottom classes.
Thus the reflection crossing has basis
\begin{equation}
 w^k x_m,\quad ew^k x_m,\qquad 0\le k\le d,
 \label{eqv:reflection-local-basis}
\end{equation}
with degrees \(a_m+k\) and \(a_m+n+k\).

For \(j\ge1\), let \(\rho_j:G\to O(2)\) multiply rotation angles
by \(j\) and preserve the chosen reflection.
The pullbacks of the two plane classes are
\begin{equation}
 w_1(\rho_j)=w,\qquad w_2(\rho_j)=(j\bmod2)U.
 \label{eqv:weight-classes}
\end{equation}
To check the second identity, restrict to \(BSO(2)\), where a
weight-\(j\) complex line multiplies the first Chern class by \(j\).
The only possible additional degree-two term is a multiple of
\(w^2\); restriction to \(BR\), where the plane has a trivial line,
excludes that term.

\begin{lemma}[Odd and even critical modules]
\label{eqv:critical-modules}
Define homogeneous polynomials \(A_i\in\kk[w,U]\) by
\begin{equation}
 A_0=1,\qquad A_1=w,\qquad
 A_i=wA_{i-1}+UA_{i-2}\quad(i\ge2).
 \label{eqv:grassmann-recurrence}
\end{equation}
If \(m\) is odd, the \(G\)-crossing is generated over
\(\kk[w,U]\) by \(x_m\), and
\begin{equation}
 A_{d+1}x_m=A_{d+2}x_m=0.
 \label{eqv:grassmann-relations}
\end{equation}
If \(m\) is even, multiplication by \(U\) is injective on the
\(G\)-crossing, and that crossing is generated over \(\kk[w,U]\)
by \(x_m\) and a chosen class \(y_m\) of degree \(a_m+n\)
whose reflection restriction is \(ex_m\).
\end{lemma}

\begin{proof}
The \(G\)-action on frames is the right action through \(\rho_m\),
up to the convention for inverse precomposition.
For odd \(m\), the map from the critical Borel space to
\(\mathrm{Gr}_2(\mathbb R^{n+1})\) has fiber the classifying space
of the cyclic group of order \(m\).
Local choices of frames give local trivializations.
That fiber has no positive-degree mod-two cohomology, by transfer,
so the map induces a cohomology isomorphism.
The tautological plane pulls back to the plane of \(\rho_m\);
its classes are \(w,U\) by \eqref{eqv:weight-classes}.

These two classes generate the Grassmann cohomology.
Indeed, projectivize the tautological plane.
On the resulting ordered-line flag space, let \(s,t\) be the first
Stiefel--Whitney classes of its two lines.
Over the Grassmannian, the projective-bundle formula gives a free
module with basis \(1,s\).
Over \(\mathbb{RP}^n\), the same flag space is the projectivization
of the first line's perpendicular bundle, so its ring is generated
by \(s,t\).
Since \(s+t=w\), \(st=U\), and \(s^2=ws+U\), every polynomial
in \(s,t\) reduces to \(P(w,U)+sQ(w,U)\).
Uniqueness of the free-module expansion implies that every class
pulled back from the Grassmannian is a polynomial in \(w,U\).
The complementary bundle has rank \(d\), with total
Stiefel--Whitney class \((1+w+U)^{-1}\).
Its components are the polynomials \(A_i\), proving
\eqref{eqv:grassmann-relations} and the odd-stage assertion.

If \(m\) is even, \(\tau\) acts trivially on the critical manifold.
Consequently
\[
 H_K^*(C_m)=H_R^*(C_m)[v].
\]
The same statement holds for the crossing modules after the Thom
shift, with the actual equivariant negative bundle retained.
The polynomial \(U=v(v+w)\) is monic in \(v\), hence acts injectively
even though the coefficient ring in
\eqref{eqv:reflection-critical-ring} has zero divisors.
The injection in \eqref{eqv:weight-two-gysin} makes \(U\) injective
on the \(G\)-crossing as well.
Sequence \eqref{eqv:standard-gysin} therefore identifies the quotient
of the \(G\)-crossing by \(U\) with the \(R\)-crossing.
Choose a lift \(y_m\) of \(ex_m\); the bottom \(x_m\) already lifts.
Their \(w\)-multiples generate modulo \(U\), by
\eqref{eqv:reflection-local-basis}.
Subtracting such a combination and dividing the remainder by \(U\)
reduces degree by two, so induction from the bottom degree proves
generation.
\end{proof}

\subsection{Collapse on all round intervals}

\begin{lemma}[Interval collapse and supported extensions]
\label{eqv:interval-collapse}
Fix \(m_0\ge1\).
For \(H=R,G\) and \(p\ge m_0\), the crossing sequence is short exact:
\begin{equation}
 0\longrightarrow H_H^j(Z_p,Z_{p-1})
 \longrightarrow H_H^j(Z_p,Z_{m_0-1})
 \longrightarrow H_H^j(Z_{p-1},Z_{m_0-1})
 \longrightarrow0.
 \label{eqv:interval-short-exact}
\end{equation}
Restriction induces an isomorphism
\begin{equation}
 H_H^j(\Lambda,Z_{m-1})
 \longrightarrow H_H^j(Z_m,Z_{m-1})
 \qquad\text{if }j<a_{m+1}.
 \label{eqv:unique-relative-lift}
\end{equation}
\end{lemma}

\begin{proof}
First use \(H=R\), inducting on \(p\).
By \eqref{eqv:reflection-local-basis}, the earlier interval module
is generated over the ambient classes \(w,e\) by chosen extensions
of its bottoms \(x_j\), \(m_0\le j<p\).
The connecting homomorphism sends such a bottom into the new crossing
in degree \(a_j+1<a_p\), where the crossing is zero.
It is linear over \(w,e\), hence vanishes on the whole earlier module.
This proves \eqref{eqv:interval-short-exact} for \(R\), and also
justifies the inductive generation assertion.

For \(G\), Lemma~\ref{eqv:critical-modules} gives the same induction
with the bottom generators and, at even stages \(j\), the chosen
top generators of degree \(a_j+n\).
Bottoms again have zero connecting image.
For a top, the only potentially relevant degree is \(a_j+n+1\).
If \(d>2\), then \(n+1=d+2<2d\), so this degree is less than \(a_p\).
If \(d=2\), equality can occur only when \(j=p-1\) is even.
In this case \(p\) is odd, and the target is its one-dimensional
bottom \(G\)-crossing.
Restriction to the corresponding reflection bottom is injective:
both bottom Thom classes restrict to the nonzero bottom class.
The reflection connecting map has already been proved zero.
Naturality therefore kills the \(G\)-connecting map in this last
case as well.  This proves \eqref{eqv:interval-short-exact} in every
dimension \(n\ge3\).

We pass to the global pair degree by degree.
Every singular chain in the Borel construction is finite and has
bounded round energy, since energy descends to a continuous invariant
function.  Thus
\[
 H_j^H(\Lambda,Z_{m-1})
 =\varinjlim_{p\ge m}H_j^H(Z_p,Z_{m-1}).
\]
Universal coefficients over the field \(\kk\) identify its dual with
the inverse limit of the degree-\(j\) cohomology groups.
The local crossing groups vanish below \(a_p\), which tends to
infinity.  Sequence \eqref{eqv:interval-short-exact} consequently
stabilizes in every degree.
For \(j<a_{m+1}\), no subsequent crossing contributes, proving
\eqref{eqv:unique-relative-lift}.
All limits in this argument are taken in a fixed degree; the total
cohomology module is the graded direct sum.
\end{proof}

Since \(Z_0\) retracts equivariantly to \(M\), the triple exact sequence
identifies \(H_H^*(Z_p,Z_0)\) with \(H_H^*(Z_p,M)\), and likewise
\(H_H^*(\Lambda,Z_0)\) with \(H_H^*(\Lambda,M)\).
Thus the first interval filtration computes precisely the global
groups relative to constants used below.

Choose the relative extension
\[
 \widetilde\zeta_m\in H_G^{a_m}(\Lambda,Z_{m-1})
\]
of the bottom Thom class.
Map it to \(H_G^{a_m}(\Lambda,M)\), and call the resulting class
\(\zeta_m\).
Let \(X_m=\res_R\zeta_m\).
For even \(m\), choose \(y_m\) as in
Lemma~\ref{eqv:critical-modules}.
Since
\begin{equation}
 n=d+1<2d,
 \label{eqv:strict-top-gap}
\end{equation}
it has a unique relative extension in degree \(a_m+n\).
Call its image in \(H_G^*(\Lambda,M)\) \(Y_m\).
Restriction and cup by \(e\) commute with the relative extensions;
both have local value \(ex_m\).
Uniqueness in \eqref{eqv:unique-relative-lift} therefore gives
\begin{equation}
 \res_RY_m=eX_m\qquad(m\ \text{even}).
 \label{eqv:even-top-restriction}
\end{equation}

Similarly, \(w^n\res_R\widetilde\zeta_m\) has local value zero
by \eqref{eqv:reflection-critical-ring}.
Its offset is \(n<2d\), so uniqueness gives
\begin{equation}
 w^nX_m=0.
 \label{eqv:global-w-relation}
\end{equation}
The products \(w^kX_m,ew^kX_m\), \(0\le k\le d\), now form a basis
of \(B_\infty\).
To see independence, restrict a finite relation to the lowest round
stage occurring in it: later classes vanish there, while the
remaining leading terms form the local basis
\eqref{eqv:reflection-local-basis}.
Repeat with successive stages.
For spanning, in each fixed degree the interval short exact sequences
give precisely the sum of the local dimensions, and only finitely
many stages contribute.
Equation \eqref{eqv:global-w-relation} and multiplication by \(e\)
show that \(w^n\) annihilates the entire relative module.

\subsection{Odd-stage torsion}

The remaining full-group relation follows from one integration
calculation.  This will also keep the dimension-three case within
the strict extension range.

\begin{lemma}[Torsion of odd bottoms]
\label{eqv:odd-torsion}
For every odd \(m\) and every \(n\ge3\),
\[
 U^n\zeta_m=0.
\]
When \(n=3\), the supported relative lift also satisfies
\[
 w^3\widetilde\zeta_m=0,\qquad
 U(w^2+U)\widetilde\zeta_m=0.
\]
\end{lemma}

\begin{proof}
The recurrence \eqref{eqv:grassmann-recurrence} implies
\begin{equation}
 A_i^2+A_{i-1}A_{i+1}=U^i\qquad(i\ge1).
 \label{eqv:determinant-identity}
\end{equation}
Indeed, substitution of the recurrence makes the left side for
\(i\ge2\) equal to \(U\) times the left side for \(i-1\);
the initial value is \(w^2+A_2=U\).

At an odd crossing, \eqref{eqv:grassmann-relations} gives
\[
 UA_dx_m=(A_{d+2}+wA_{d+1})x_m=0.
\]
The class \(A_dx_m\) is nonzero: on restriction to \(R\), it becomes
\(w^dx_m\), which is nonzero in
\eqref{eqv:reflection-local-basis}.
The relevant part of the standard-plane Gysin sequence is
\[
 H_R^{a_m+d+1}(Z_m,Z_{m-1})
 \xrightarrow{\int} H_G^{a_m+d}(Z_m,Z_{m-1})
 \xrightarrow{U} H_G^{a_m+d+2}(Z_m,Z_{m-1}).
\]
Its source is one-dimensional, spanned by \(ex_m\).
Since its image contains the nonzero class \(A_dx_m\), exactness
over \(\kk\) forces
\begin{equation}
 \int ex_m=A_dx_m.
 \label{eqv:local-integration}
\end{equation}

Apply the same integration map to
\(e\,\res_R\widetilde\zeta_m\) in the pair
\((\Lambda,Z_{m-1})\).
Naturality and \eqref{eqv:local-integration} identify its restriction
to the crossing with \(A_dx_m\).
Both sides have offset \(d<2d\), so
\eqref{eqv:unique-relative-lift} gives
\[
 \int\bigl(e\,\res_R\widetilde\zeta_m\bigr)
   =A_d\widetilde\zeta_m.
\]
The next map in Gysin is multiplication by \(U\), hence
\begin{equation}
 UA_d\widetilde\zeta_m=0.
 \label{eqv:integrated-relation}
\end{equation}
Also \(A_{d+1}\widetilde\zeta_m=0\), because its local value is zero
and \(d+1<2d\).
Multiplying \eqref{eqv:determinant-identity} at \(i=d\) by
\(U\widetilde\zeta_m\) therefore gives
\[
 U^{d+1}\widetilde\zeta_m
 =UA_d^2\widetilde\zeta_m+
   UA_{d-1}A_{d+1}\widetilde\zeta_m=0.
 \]
Its image in \(H_G^*(\Lambda,M)\) is \(U^n\zeta_m=0\).

For \(n=3\), one has \(d=2\), \(A_2=w^2+U\), and \(A_3=w^3\).
The two intermediate global relations are exactly
\[
 w^3\widetilde\zeta_m=0,\qquad
 U(w^2+U)\widetilde\zeta_m=0.
\]
Their proof used uniqueness only at offsets two and three, both
strictly less than the next-bottom gap four.
It did not extend an offset-four relation across that gap.
\end{proof}

The basis construction, \eqref{eqv:even-top-restriction}, and
Lemma~\ref{eqv:odd-torsion} complete the proof of
Theorem~\ref{thm:global-generators}.
For later use, the degree intervals for the bottom family are
\([a_m,a_m+d]\), and those for the top family are
\([a_m+n,a_m+n+d]\).
Intervals within either family are disjoint, since their starting
degrees differ by \(2d\); a bottom and a top can have the same degree.
Thus restrictions of tops will be counted modulo all restricted
bottoms.
The next sections compare those global restrictions with the extra
classes present at a finite cut.
